% =================================================
% Farben
% =================================================
\definecolor{codeblue}{HTML}{4178f2}
\definecolor{lieblingsfarbe}{cmyk}{79,0.11,1,0.25}

% Hyperref-Optionen (Beamer lädt hyperref selbst)
\hypersetup{
	hidelinks,
	colorlinks=true,
	allcolors=codeblue
}

% =================================================
% Schrift & Kodierung
% =================================================
\usepackage{lmodern}
\usepackage{inconsolata}

% Entfernt: \newunicodechar{̈}{\"} (führt zu Endlosschleife!)

% =================================================
% Mathematik
% =================================================
\usepackage{amsmath}      % von Beamer unterstützt
\usepackage{amssymb}
\usepackage{amsfonts}
\usepackage[centerdots]{mathtools}
\usepackage{dsfont}
\usepackage{stmaryrd}
\usepackage{nicefrac}

% =================================================
% Grafiken & TikZ
% =================================================
\usepackage{graphicx}
\usepackage{tikz}
\usepackage{pgfplots}
\pgfplotsset{compat=1.18}

\usetikzlibrary{
	chains,
	shapes.multipart,
	shapes,
	arrows,
	arrows.meta,
	positioning,
	fit,
	shapes.geometric
}

% TikZ-Stile
\tikzset{
	heapnode/.style={draw, circle, inner sep=0pt, minimum size=6mm},
	heaparrow/.style={-latex, thick},
	long node/.style={draw, rounded corners, minimum width=7.5cm, minimum height=1cm},
	short node/.style={draw, rounded corners, minimum width=2.5cm, minimum height=1cm},
	square node/.style={draw, rounded corners, minimum width=2.5cm, minimum height=2.5cm},
	label node/.style={font=\small, align=center}
}

% =================================================
% Tabellen & Layout
% =================================================
\usepackage{array}
\usepackage{booktabs}
\usepackage{float}
\usepackage{colortbl}

% =================================================
% Aufzählungen
% =================================================
\usepackage{enumitem}
\setlist[itemize]{label=$\bullet$, leftmargin=*, nosep}
\setlist[enumerate]{label=\arabic*., leftmargin=*, nosep}

% =================================================
% Listings / Code
% =================================================
\usepackage{listings}

\lstdefinelanguage{pseudocode}{
	keywords={struct,null,function,error,nil,if,then,else,for,while,do,in,True,False,Array,set,repeat,break,continue,Eingabe,Ausgabe,return,print,exit},
	sensitive=false,
	delim=[l][keywordstyle]{:},
	comment=[l][commentstyle]{\#}
}

\lstset{
	frame=single,
	framesep=0mm,
	framexleftmargin=7mm,
	xleftmargin=8mm,
	framerule=1.5pt,
	rulecolor=\color{black},
	backgroundcolor=\color{white},
	basicstyle=\ttfamily,
	keywordstyle=\bfseries\color{codeblue},
	commentstyle=\color{gray},
	numbers=left,
	stepnumber=1,
	numbersep=1mm,
	numberstyle=\sffamily\color{gray!80!black}\footnotesize,
	numberblanklines=true,
	escapeinside={\%*}{*)},
	inputencoding=utf8
}

% =================================================
% Algorithmen
% =================================================
\usepackage{algorithm}
\usepackage{algpseudocode}

% =================================================
% Sonstiges
% =================================================
\usepackage{textcomp}
\usepackage{tipa}
\usepackage{blindtext}
\usepackage{tcolorbox}
\usepackage{csquotes}

% =================================================
% Beamer-Theoremumgebungen
% =================================================
\newtheorem{satz}{Satz}
\newtheorem{fundamentalsatz}{Fundamentalsatz}
\newtheorem{intuition}{Intuition}
\newtheorem{intuitionUndBeispiel}{Intuition und Beispiel}
\newtheorem{beispiel}{Beispiel}
\newtheorem{korollar}{Korollar}
\newtheorem{bemerkung}{Bemerkung}
\newtheorem{vorbemerkung}{Vorbemerkung}

% Beweisumgebung
\newenvironment{beweis}[1][\textbf{Beweis}]{%
	\begin{proof}[#1]
	}{%
	\end{proof}
}

% =================================================
% Trennfolie
% =================================================
\newcommand{\trennfolie}[1]{
	\begin{frame}[plain]
		\centering
		\vfill
		\Huge \textbf{#1}
		\vfill
	\end{frame}
}

% =================================================
% Metadaten
% =================================================
\author{Yavuzâlp Dal}
\title{Einf\"uhrung in die Kerndichtesch\"atzung}
\institute{Angewandte Statistik}
\date{\today}
